%MB9T5










\setcounter{paragraph}{0}
\setcounter{ex}{0} 
\PhanI
\begin{ex} %[3P4Y4-2][Loai3] 
\nguon{QG - 2020} Một sóng điện từ có tần số 75 kHz đang lan truyền trong chân không. Lấy $c=3.10^8\ m/s$. Sóng này có bước sóng là
\choice
{2000 m} 
{0,5 m} 
{\True 4000 m}
{0,25 m}
\loigiai{$\lambda=\dfrac{c}{f}=4000\ m$}
%<MyLT>
\end{ex} \haidong{20}

\begin{ex} %[3P4Y4-2][Loai3]
\nguon{QG - 2017} Một sóng điện từ có tần số 90 MHz, truyền trong không khí với tốc độ ${3.10^8}$ m/s thì có bước sóng là
\choice
{\True 3,333 m}
{3,333 km}
{33,33 m}
{33,33 km}
\loigiai{
Bước sóng của sóng $\lambda =\dfrac{c}{f}=\dfrac{3.10^8}{90.10^6}=3,333\ m.$
}
%<MyLT2>
\end{ex} \haidong{20}


\begin{ex}%CTST
Một vệ tinh nhân tạo chuyển động ở độ cao 575 km so với mặt đất phát sóng vô tuyến có tần số 92,4 MHz với công suất 25 kW về phía mặt đất. Bỏ qua sự hấp thụ sóng của khí quyển. Cường độ sóng nhận được bởi một máy thu vô tuyến ở mặt đất ngay dưới vệ tinh là
\choice
{$6,02.10^{-4}\ W/m^2$}
{\True $6,02.10^{-9}\ W/m^2$}
{$6,02.10^{-6}\ W/m^2$}
{$6,02.10^{-3}\ W/m^2$}
\loigiai{
\begin{itemize}
\item Cường độ sóng: $I=\dfrac{P}{4\pi r^2}=\dfrac{25.10^3}{4\pi.(575.10^3)^2}=6,02.10^{-9}\ W/m^2$
\end{itemize}
}
\end{ex} \haidong{20}
\begin{ex} %[3P4Y3-1][Loai1] 
\nguon{QG - 2018} Khi nói về sóng điện từ, phát biểu nào sau đây \textbf{sai}?
\choice
{Sóng điện từ là sóng ngang}
{Sóng điện từ mang năng lượng}
{\True Sóng điện từ không truyền được trong chân không}
{Sóng điện từ có thể phản xạ, khúc xạ hoặc giao thoa}
\loigiai{
Sóng điện từ truyền được trong chân không. Nên C sai}
\end{ex} \haidong{20}
\begin{ex} %[3P4B3-1][Loai1] 
Khi sóng âm và sóng điện từ truyền từ không khí vào nước thì
\choice
{tốc độ truyền sóng âm và bước sóng của sóng điện từ đều giảm}
{tốc độ truyền sóng âm giảm, bước sóng của sóng điện từ tăng}
{\True bước sóng của sóng âm tăng, bước sóng của sóng điện từ giảm}
{bước sóng của sóng âm và bước sóng của sóng điện từ đều tăng}
\loigiai{
Khi sóng âm và sóng điện từ truyền từ không khí vào nước thì cả hai sóng đều có tần số không đổi
\begin{itemize}
\item Vận tốc truyền sóng của sóng âm trong nước lớn hơn trong không khí do vậy bước sóng của sóng âm tăng.
\item Chiết suất của nước đối với sóng điện từ lớn hơn chiết suất của không khí đối với sóng điện từ do vậy vận tốc truyền sóng giảm kết quả là bước sóng cũng giảm theo.
\end{itemize}
}
\end{ex} \haidong{20}


\begin{ex} %[3P4Y4-1][Loai3]
Truyền hình số vệ tinh K+ sử dụng vệ tinh Vinasat. Sóng vô tuyến truyền hình K+ thuộc dải
\choice
{sóng trung}
{sóng ngắn}
{\True sóng cực ngắn}
{sóng dài}
\loigiai{Truyền hình vệ tinh sử dụng sóng cực ngắn.
}
\end{ex} \haidong{20}

\begin{ex} %[3P4B4-1][Loai12]
Trong "máy bắn tốc độ" xe cộ trên đường
\choice
{\True có cả máy phát và máy thu sóng vô tuyến}
{chỉ có máy thu sóng vô tuyến}
{chỉ có máy phát sóng vô tuyến}
{không có máy phát và máy thu sóng vô tuyến}
\loigiai{
Máy bắn tốc độ gồm cả máy phát và máy thu sóng vô tuyến.
}
\end{ex} \haidong{20}


























\begin{ex} 
\nguon{QG - 2019} Trong chân không bức xạ có bước sóng nào sau đây là bức xạ hồng ngoại
\choice
{\True 900 nm}
{250 nm}
{450 nm}
{600 nm}
\loigiai{
Bức xạ hồng ngoại có bước sóng dài hơn 760 nm.
}
\end{ex} \haidong{20}
\begin{ex}%[3P5Y8-1][Loai1]
Tia hồng ngoại
\choice
{không truyền được trong chân không}
{là ánh sáng nhìn thấy, có màu hồng}
{không phải là sóng điện từ}
{\True được ứng dụng để sưởi ấm}
\loigiai{
Tia hồng ngoại có tác dụng nhiệt rất mạnh nên được ứng dụng để sưởi ấm.
}
\end{ex} \haidong{20}
\begin{ex}%[3P5Y8-1][Loai1]
Tia tử ngoại \textbf{không} có tác dụng nào sau đây?
\choice
{Quang điện}
{Kích thích phát quang}
{\True Chiếu sáng}
{Sinh lí}
\loigiai{
Tia tử ngoại là bức xạ không nhìn thấy nên không có tác dụng thắp sáng.
}
\end{ex} \haidong{20}
\begin{ex} %[3P5Y8-1][Loai1]
\nguon{QG - 2020} Phát biểu nào sau đây là \textbf{sai}?
\choice
{Tia X có tác dụng sinh lý}
{Tia X làm ion hóa trong không khí}
{Tia X có bước sóng nhỏ hơn bước sóng của ánh sáng tím}
{\True Tia X có bước sóng lớn hơn bước sóng của tia hồng ngoại}
\loigiai{Tia X có bước sóng nhỏ hơn bước sóng của tia hồng ngoại}
\end{ex} \haidong{20}
\begin{ex}%[3P5Y8-1][Loai1]
\nguon{QG - 2015} Khi nói về tia X, phát biểu nào sau đây là đúng?
\choice
{\True Tia X có tác dụng sinh lí: nó hủy diệt tế bào}
{Tia X có tần số nhỏ hơn tần số của tia hồng ngoại}
{Tia X có bước sóng lớn hơn bước sóng của ánh sáng nhìn thấy}
{Tia X có khả năng đâm xuyên kém hơn tia hồng ngoại}
\loigiai{
Tia X có tác dụng sinh lí: nó hủy diệt tế bào.
}
\end{ex} \haidong{20}
\begin{ex}%[3P5Y8-1][Loai1]
Để kiểm tra hành lí khách đi máy bay, người ta dùng 
\choice
{tia đỏ}
{tia tử ngoại}
{\True tia X}
{tia hồng ngoại}
\loigiai{
Dựa vào tính đâm xuyên của tia X người ta kiểm tra hành lý.
}
\end{ex} \haidong{20}
\begin{ex}%[3P5Y8-1][Loai1]
Trong chân không, xét các tia: tia hồng ngoại, tia tử ngoại, tia X và tia đơn sắc lục. Tia có bước sóng nhỏ nhất là
\choice
{tia hồng ngoại}
{tia đơn sắc lục}
{\True tia X}
{tia tử ngoại}
\loigiai{
Ta có thứ tự các bức xạ có bước sóng tăng dần trong chân không: Tia gamma $(\gamma ) $ , tia Rơn–ghen (X), tia tử ngoại, ánh sáng tím, ánh sáng đỏ, tia hồng ngoại, sóng vô tuyến.
}
\end{ex} \haidong{20}
\begin{ex}%[3P5Y8-1][Loai1]
Bốn bức xạ: ánh sáng nhìn thấy, tia hồng ngoại, tia X và tia $\gamma$. Các bức xạ này được sắp xếp theo thứ tự bước sóng tăng dần là
\choice
{tia X, ánh sáng nhìn thấy, tia $\gamma $, tia hồng ngoại}
{tia $\gamma $, tia X, tia hồng ngoại, ánh sáng nhìn thấy}
{\True tia $\gamma $, tia X, ánh sáng nhìn thấy, tia hồng ngoại}
{tia $\gamma $, ánh sáng nhìn thấy, tia X, tia hồng ngoại}
\loigiai{
\begin{itemize}
\item Tia gamma ( $\gamma $) có cùng bản chất với tia X nhưng có bước sóng ngắn hơn.
\item Áp dụng cả câu 1 vậy ta có thứ tự bước sóng tăng dần: tia gamma, tia X, anh sáng nhìn thấy, tia hồng ngoại.
\end{itemize}
}
\end{ex} \haidong{20}
\begin{ex}%[3P5Y8-1][Loai1]
Bước sóng của tia hồng ngoại nhỏ hơn bước sóng của
\choice
{\True sóng vô tuyến}
{tia Rơn-ghen}
{ánh sáng tím}
{ánh sáng đỏ}
\loigiai{
Ta có thứ tự các bức xạ có bước sóng tăng dần trong chân không: Tia gamma $(\gamma )$, tia Rơn–ghen (X), tia tử ngoại, ánh sáng tím, ánh sáng đỏ, tia hồng ngoại, sóng vô tuyến.
}
\end{ex} \haidong{20}
\begin{ex}%[3P5Y8-1][Loai1]
Cho các tia: hồng ngoại, tử ngoại, đơn sắc màu lục và Rơn-ghen. Trong cùng một môi trường truyền, tia có bước sóng dài nhất là
\choice
{\True tia hồng ngoại}
{tia tử ngoại}
{tia đơn sắc màu lục}
{tia Rơn-ghen}
\loigiai{
Tia hồng ngoại có bước sóng dài nhất.
}
\end{ex} \haidong{20}
\begin{ex}%[3P5Y8-1][Loai1]
Cho các nguồn phát bức xạ điện từ chủ yếu (xem mỗi dụng cụ phát một bức xạ) gồm: Bàn là áo quần (I), đèn quảng cáo (II), máy chụp kiểm tra tổn thương xương ở cơ thể người (III), điện thoại di động (IV). Các bức xạ do các nguồn trên phát ra sắp xếp theo thứ tự tần số giảm dần là
\choice
{IV, I, III, II}
{IV, II, I, III}
{III, IV, I, II}
{\True III, II, I, IV}
\loigiai{
\begin{itemize}
\item IV phát ra sóng vô tuyến; nhiệt độ từ I chỉ đủ để phát ra tia hồng ngoại; II có nhiệt độ lớn hơn I, III phát ra tia X.
\item Vậy có thứ tự tần số giảm dần là III, II, I, IV.
\end{itemize}
}
\end{ex} \haidong{20}

\setcounter{paragraph}{0}
\setcounter{ex}{0} 
\PhanII
\begin{ex}
Tia tử ngoại (UV) là bức xạ điện từ có năng lượng cao hơn ánh sáng tím và mang lại nhiều ứng dụng cũng như tác hại trong đời sống.
\choiceTF[t]
{\True Vật có nhiệt độ trên $3000^\circ C$ phát ra tia tử ngoại rất mạnh}
{Tia tử ngoại không bị thủy tinh hấp thụ}
{\True Tia tử ngoại có tác dụng sinh học: diệt vi khuẩn, hủy diệt tế bào da}
{\True Tia tử ngoại tác dụng lên phim ảnh}
\loigiai{
\begin{itemchoice}
\itemch
\itemch
\itemch
\itemch
\end{itemchoice}
}
\end{ex} \haidong{20}
\begin{ex}
Tia X là bức xạ điện từ có bước sóng rất ngắn và năng lượng lớn, được dùng rộng rãi trong y học cũng như nghiên cứu vật chất.
\choiceTF[t]
{\True Tia X làm ion hóa không khí}
{Tia X có bước sóng lớn hơn bước sóng của tia tử ngoại}
{\True Tia X có khả năng đâm xuyên mạnh}
{\True Tia X làm phát quang một số chất}
\loigiai{
\begin{itemchoice}
\itemch
\itemch
\itemch
\itemch
\end{itemchoice}
}
\end{ex} \haidong{20}
\begin{ex}
Trong kĩ thuật và y – sinh, tia X được tạo ra khi chùm electron năng lượng cao va vào bia kim loại nặng và tia X có tính đâm xuyên mạnh.
\choiceTF[t]
{\True Ta có thể tạo ra tia X nhờ ống tia X: chùm electron có vận tốc lớn đập vào đối catot làm bằng kim loại có nguyên tử lượng lớn như Platin, làm bật ra chùm tia X}
{\True Tia X có cùng bản chất với tia hồng ngoại}
{Tia X có vận tốc lớn hơn vận tốc ánh sáng}
{Tia X được phát ra từ các vật nóng sáng trên $500^\circ C$}
\loigiai{
\begin{itemchoice}
\itemch
\itemch
\itemch
\itemch
\end{itemchoice}
}
\end{ex} \haidong{20}
\begin{ex}
Trong chẩn đoán hình ảnh và nghiên cứu vật chất thì tia X nổi bật nhờ năng lượng rất lớn, khả năng đâm xuyên mạnh và tác dụng ion hoá.
\choiceTF[t]
{\True Tia X là các bức xạ điện từ có bước sóng từ $10^{-11} \ m$ đến $10^{-8} \ m$}
{Tia X không có trong ánh sáng của Mặt Trời khi truyền đến Trái Đất}
{\True Tia X được phát ra từ bản kim loại nặng, khó nóng chảy khi có một chùm electron có động năng lớn đập vào}
{\True Tia X có tác dụng mạnh lên phim ảnh, làm ion hóa không khí}
\loigiai{
\begin{itemchoice}
\itemch
\itemch
\itemch
\itemch
\end{itemchoice}
}
\end{ex} \haidong{20}

\setcounter{ex}{0}
\PhanIII
\begin{ex} %[3P4Y4-2][Loai3] 
\nguon{QG - 2017} Một sóng điện từ có tần số 30 MHz truyền trong chân không với tốc độ $3.10^8$ m/s thì có bước sóng là
\shortans[oly]{10}
\loigiai{
Bước sóng của sóng điện từ: $\lambda =\dfrac{c}{f}=\dfrac{3.10^8}{30.10^6}=10\ m$.
}
%<MyLT>
\end{ex} \haidong{20}
\begin{ex}
Trong chân không, tốc độ truyền sóng điện từ bằng $x.10^8\ m/s$. Tìm x.
\shortans[oly]{3}
\loigiai{}
\end{ex}
\begin{ex} %[3P4Y4-2][Loai3] 
\nguon{ĐH - 2013} Sóng điện từ có tần số 10 MHz truyền trong chân không với bước sóng là bao nhiêu mét?
\shortans[oly]{30}
\loigiai{
Ta có: $\lambda=\dfrac{c}{f}=30\ m$.
}%<MyLT1>
\end{ex} \haidong{20}
\loai{Tính T, f}
\begin{ex} %[3P4Y4-2][Loai4]
Sóng FM của đài Hà Nội có bước sóng $\lambda =\dfrac{10}{3}\ m$. Tần số f là bao nhiêu MHz?
\shortans[oly]{90}
\loigiai{
Sóng FM của đài Hà Nội là sóng điện từ lan truyền trong không gian với vận tốc $c = 3.10^8\ m/s$
$\Rightarrow$ Tần số $f = \dfrac{v}{\lambda} =\dfrac{3.10^8}{\dfrac{10}{3}}=90.10^6\ Hz=90\ MHz$. 
}
%<MyLT>
\end{ex} \haidong{20}
\begin{ex}%[3P4Y4-2][Loai4]
\nguon{QG - 2019} Một sóng điện từ lan truyền trong chân không có bước sóng 6000 m. Lấy $c=3.10^8$ m/s. Biết trong sóng điện từ, thành phần điện trường tại một điểm biến thiên điều hòa với chu kì T. Giá trị của T là $x.10^{-5}\ s$. Tìm x.
\shortans[oly]{2}
\loigiai{
Ta có: $\lambda =c.T\Rightarrow T=\dfrac{\lambda}{c}=\dfrac{6000}{3.10^8}=2.10^{-5}\ s.$}
%<MyLT2>
\end{ex} \haidong{20}
\begin{ex}%[3P4Y4-2][Loai4]
\nguon{QG - 2019} Một sóng điện từ lan truyền trong chân không có bước sóng 3000 m. Lấy $c = 3.10^8$ m/s. Biết trong sóng điện từ, thành phần điện trường tại một điểm biến thiên với tần số f. Giá trị của f là $x.10^5\ Hz$. Tìm x.
\shortans[oly]{2}
\loigiai{
Ta có: $\lambda =\dfrac{c}{f}\Rightarrow f=\dfrac{c}{\lambda}=\dfrac{3.10^8}{3000}=10^5\ Hz$}
\end{ex} \haidong{20}
